\documentclass[11pt,a4paper]{article}
\usepackage[T1]{fontenc}
\usepackage[utf8]{inputenc}
\usepackage{lmodern,microtype}
\usepackage[margin=26mm]{geometry}
\usepackage{amsmath,amssymb,amsthm,mathtools}
\usepackage{booktabs,tabularx,longtable,array}
\usepackage{graphicx,xcolor}
\usepackage{listings}
\usepackage[colorlinks=true,linkcolor=blue!45!black,citecolor=blue!45!black,urlcolor=blue!45!black]{hyperref}
\usepackage{fancyhdr}
\pagestyle{fancy}\fancyhf{}
\fancyhead[L]{\small Exact acceleration of \texttt{fastunknot}}
\fancyhead[R]{\small Implementation report, v0.2.0}
\fancyfoot[C]{\thepage}
\setlength{\headheight}{14pt}
\setlength{\parskip}{4pt}
\setlength{\parindent}{15pt}
\setlength{\emergencystretch}{2em}
\newtheorem{theorem}{Theorem}[section]
\newtheorem{proposition}[theorem]{Proposition}
\newtheorem{lemma}[theorem]{Lemma}
\newtheorem{corollary}[theorem]{Corollary}
\theoremstyle{definition}\newtheorem{remark}[theorem]{Remark}
\newcommand{\F}{\mathbb F_2}
\newcommand{\Kh}{\operatorname{Kh}}
\newcommand{\rKh}{\widetilde{\operatorname{Kh}}}
\newcommand{\rank}{\operatorname{rank}}
\newcommand{\poly}{\operatorname{poly}}
\newcommand{\id}{\operatorname{id}}
\newcommand{\Unknot}{\textsc{Unknot}}
\newcommand{\Knotted}{\textsc{Knotted}}
\newcommand{\Unknown}{\textsc{Unknown}}
\lstset{basicstyle=\small\ttfamily,breaklines=true,columns=fullflexible,
  keepspaces=true,showstringspaces=false,frame=single,framerule=.3pt,
  xleftmargin=3pt,xrightmargin=3pt,aboveskip=9pt,belowskip=8pt}
\title{Exact acceleration of the supplied unknot recognizer\\[4pt]
\large Modular frontiers, compressed connected sums, and sparse cancellation}
\author{Implementation and mathematical report for \texttt{Knots.zip}}
\date{18 September 2026 \quad---\quad Version 0.2.0}

\begin{document}
\maketitle
\begin{abstract}
We replace the supplied \texttt{fast/} implementation by a compatible,
standard-library Python recognizer with three substantive improvements: an exact
one-sided modular Jones obstruction that aggregates partial smoothing states;
a validated decomposition along visible connected-sum circles, with tensor products
of reduced Khovanov homology kept compressed; and a faster Bar--Natan scanning
backend using characteristic-two algebra, bounded memoization, and fill-in-aware
cancellation. The original implementation is retained unchanged for comparisons.

On the same machine, cold-process median recognition times decrease from
35.09\,ms to 0.834\,ms for the Conway fixture, and from 1.372\,s to
1.455\,ms for two connected Conway copies. The original raw scanner does not
complete the supplied 36-crossing stress case within our 30-second budget;
the improved raw scanner completes in 1.487\,s. Some already easy inputs become
slower; all measured comparisons are reported. Independent small-cube and
state-sum tests, differential checks, evidence replay, and a large-diagram
Reidemeister invariance check support the implementation.

No global $n^{O(\log n)}$ running-time theorem is claimed. The general worst-case
bound remains exponential. There is, however, a provable exponential improvement
on a fixed-block connected-sum family, and a quasipolynomial bound on the restricted
class whose extracted blocks have $O(\log^2 n)$ crossings. We explain precisely
why small scan width alone does not control surviving Khovanov multiplicities,
and distinguish the supplied Lackenby hierarchy results from a complete executable
realization of the announced quasipolynomial algorithm.
\end{abstract}

\clearpage
\begingroup\small\setlength{\parskip}{0pt}
\tableofcontents
\endgroup
\clearpage

\section{Scope, deliverables, and the baseline}
\label{sec:scope}
The input is a \emph{classical one-component knot diagram} with $n$ crossings.
The implementation accepts planar-diagram (PD) codes, braid closures, and grid
diagrams. A PD row $(a,b,c,d)$ gives the four ports in cyclic order; the
under-strand joins ports 0 and 2, and the over-strand joins ports 1 and 3.
Each edge label must occur exactly twice. The empty PD denotes one crossing-free
circle, not the empty link. The validator checks strand connectivity and that the
specified rotation system is spherical. Links and virtual diagrams are rejected.
The crossing number $n$ throughout this report is the number of crossings in the
\emph{given diagram}, not the minimum crossing number of its knot type.

The supplied baseline first performs crossing-decreasing Reidemeister-I/II moves,
then a descending-diagram test, then computes the full Alexander polynomial. If
none decides the answer, it computes unreduced Khovanov homology over $\F$ by
Bar--Natan scanning \cite{BNfast,BNtangles}. The total rank divided by two is
used as reduced rank. Rank one detects the unknot, as explained in
Section~\ref{sec:rank}. Resource exhaustion gives \Unknown, not a knot verdict.

The new default pipeline is
\[
\begin{gathered}
\text{R1/R2}\ \longrightarrow\ \text{descending}\ \longrightarrow\ \text{visible sums}\\[3pt]
\longrightarrow\ \text{bounded Jones screen}\ \longrightarrow\ \Delta_K\ \longrightarrow\ \text{improved scan}.
\end{gathered}
\]
For a detected sum, the pipeline is applied to its factors, and stops as soon as a
nontrivial factor is certified. The decomposition itself is performed iteratively;
there is no recursion through an arbitrarily deep decomposition tree.

The package contains executable code, not calls to hypothetical hierarchy oracles.
The implementation directory is \path{fast/fastunknot}; the unchanged supplied
Python sources are in \path{baseline/fastunknot}. New modules are
\path{jones.py}, \path{factor.py}, and \path{verify.py}. The existing scanning,
preprocessing, Alexander, and command-line modules are modified. The tests,
benchmark scripts, original stress inputs, raw measurement records, this source,
and its compiled PDF are included. No external knot library, CAS, native compiler,
or network connection is needed to run the recognizer.

\begin{center}
\small
\begin{tabularx}{\linewidth}{@{}p{.34\linewidth}X@{}}
\toprule
Claim & Status in this deliverable\\
\midrule
Exact complete mathematical decision algorithm & Yes, with unbounded resources and the standard Khovanov theorems.\\
Faster practical implementation & Measured on the supplied examples and structured families; regressions are included.\\
Exponential gain on a family & Yes: visible connected sums of a fixed nontrivial Alexander-trivial knot.\\
Global quasipolynomial bound & Not established and not advertised by the API.\\
Implementation of Lackenby's hierarchy algorithm & No; its relevant ingredients and remaining implementation obligations are analyzed.\\
Formal verification of the implementation & No. Tests and replayable evidence are provided, with a stated trust boundary.\\
\bottomrule
\end{tabularx}
\end{center}

\section{The quasipolynomial target and the supplied sources}
\label{sec:lackenby}
\subsection{What the talk asserts}
The supplied talk is dated February 2021, although the accompanying preprint is
from July 2026. Physical slide 14 states the quasipolynomial recognition claim.
Slides 66--74 discuss a bound of the form $Lg^L$, starting with linear $L$ and
exponential $g$. The advertised improvements use surfaces of controlled
complexity, efficient hierarchy encodings, multisurfaces, and Heegaard splittings
with Cheeger regions. Slide 74 gives the intended logarithmic hierarchy depth;
the final slide, physical page 109, is the full control-flow diagram
\cite{Ltalk}. These are substantially different objects and operations from
Khovanov scanning.

\subsection{What the 2026 preprint contributes}
The supplied source \cite{L2026} is an important constructive foundation, not
merely a statement that an algorithm exists. Its Section 9 analyzes the number
of passes through steps of the hierarchy algorithm using the maximal hierarchy
length $L$ and maximal surface pattern-complexity $g$. Proposition 9.1 gives
\begin{equation}
N_{\rm steps}\le L(g+1)^L.
\label{eq:hierarchy}
\end{equation}
That section explicitly distinguishes the iteration count from a complete
running-time analysis of all subroutines.

It would be inaccurate to stop reading there and assert that the paper has no
bounds on hierarchy size or no compressed cutting procedure. Section 10 refines
the construction using interior parallelity bundles and obtains a \emph{linear}
hierarchy-length bound in the initial $q$-complexity. Section 11 controls uniform
handle types. Proposition 14.2 gives a polynomial-time construction of an
admissible handle structure after the relevant cut-and-simplify operations, as a
function of the number of handles and the logarithm of the surface's extended
weight, under its stated hypotheses, including an essential boundary pattern.
Theorem 14.4 supplies compressed cutting data and descriptions of the parallelity
bundle. These results must be credited in any account of the available machinery.

Their hypotheses, encoding conventions, and branch invariants still have to be
implemented and connected to the multisurface/Cheeger-region speedups in the
talk. In particular, a linear bound on hierarchy length cannot simply be inserted
into \eqref{eq:hierarchy} and called quasipolynomial. Nor can an encoded surface
of polynomial bit length safely be expanded into exponentially many normal discs.
This report does not dispute the announced target; it does not supply that
complete realization.

\subsection{The exact accounting that would suffice}
\begin{lemma}[Conditional complexity calculation]
Suppose an implemented hierarchy algorithm obeys \eqref{eq:hierarchy}, has
$L\le b\log(n+2)$ and $g\le(n+2)^a$, and every operation, including all data
manipulation, costs at most $(n+2)^c$ bit operations. Here $a,b,c$ are constants
independent of the input. Its total running time is $n^{O(\log n)}$.
\end{lemma}
\begin{proof}
The logarithm of the product of the step count and step cost is at most
\[
\log L+L\log(g+1)+c\log(n+2)=O(\log^2(n+2)).
\]
Exponentiation gives the result. All encoded integer sizes and expanded output
sizes must be included in the per-step hypothesis.
\end{proof}

For completeness, the base-$(g+1)$ potential underlying \eqref{eq:hierarchy} is
\[
\Phi=\sum_{i=1}^{L} c_i(g+1)^{L-i},\qquad 0\le c_i\le g.
\]
Decreasing digit $c_j$ by one while replacing all following digits by arbitrary
values at most $g$ decreases $\Phi$ by at least
\[
(g+1)^{L-j}-g\sum_{t=0}^{L-j-1}(g+1)^t=1.
\]
This controls the number of simplification cycles. It does not control the cost
of finding the next surface or of representing the resulting manifold.

A faithful implementation project therefore needs concrete, tested operations on
uniform handle structures and boundary patterns; compressed normal-surface and
parallelity-bundle data; compression and weak-reduction certificates; and a
verified implementation of the logarithmic-depth strategy. None is replaced here
by enumeration hidden behind a promising function name. The improvements below
instead work directly with the supplied executable algorithm.

\section{An exact modular Jones obstruction}
\label{sec:jones}
\subsection{Normalization and the one-sided decision rule}
Put $\delta=-A^2-A^{-2}$. For a nonempty PD, let $s\in\{0,1\}^n$ select a
smoothing at each crossing, with smoothing 0 pairing $(a,b),(c,d)$ and smoothing
1 pairing $(a,d),(b,c)$. Write $c(s)$ for the number of closed circles. We use
the unnormalized bracket
\begin{equation}
B_D(A)=\sum_s A^{n-2|s|}\delta^{c(s)}.
\label{eq:bracket}
\end{equation}
Thus a crossing-free circle has bracket $\delta$. With the oriented crossing-sign
convention used consistently by the code, writhe normalization gives
\begin{equation}
D\text{ represents the unknot}\quad\Longrightarrow\quad
B_D(A)=\delta(-A^3)^{w(D)}.
\label{eq:unknotbracket}
\end{equation}
This is the Kauffman state model with a loop-normalization factor
\cite{Kauffman}. The usual convention with bracket one on a circle differs by
that factor.

Choose an integer modulus $m\ge2$ and an integer $a$ coprime to $m$. Evaluation
at $A=a$ defines a homomorphism from $\mathbb Z[A,A^{-1}]$ to
$\mathbb Z/m\mathbb Z$. In particular, negative writhe causes no problem:
$-a^3$ is a unit. The default uses $a=2$ and $m=1\,000\,000\,007$.

\begin{proposition}[Sound obstruction]
If the two sides of \eqref{eq:unknotbracket}, evaluated modulo $m$, differ, then
$D$ represents a nontrivial knot. Equality is inconclusive.
\end{proposition}
\begin{proof}
An equality in the Laurent polynomial ring remains an equality under every ring
homomorphism. Therefore a modular inequality contradicts the necessary unknot
identity. The converse does not follow: evaluation can lose information, even
when two Laurent polynomials differ.
\end{proof}

There is no probabilistic acceptance, no primality assumption on $m$, and no
modular division by $\delta$. Even $\delta=0$ is safe, though the test then has no
power on a nonempty knot diagram. The implementation never turns agreement of
one evaluation, or even agreement of the full Jones polynomial, into an unknot
verdict.

\subsection{Frontier aggregation instead of the full state cube}
After some crossings have been processed, the smoothing in that region consists
of closed circles and paths joining the current boundary edge midpoints. Absorb
the circles into the coefficient. The remaining combinatorial state is a perfect
matching $\pi$ of the boundary labels. Maintain a dictionary
\[
F_i:\{\text{boundary matchings}\}\longrightarrow\mathbb Z/m\mathbb Z.
\]
Its value at $\pi$ is the sum of the weights of \emph{all} partial smoothing
choices inducing $\pi$.

To add a crossing, try its two smoothings for each dictionary entry. A small
disjoint-set forest joins old matching arcs and new smoothing arcs. Every
connected component has either two new boundary ends or no boundary ends. The
former gives a new matching arc; the latter contributes a factor $\delta$.
Multiply by $a$ or $a^{-1}$, and add the result to the new matching's residue.
Entries that cancel to zero are deleted immediately. No cobordism compositions
or Khovanov basis vectors are used by this routine.

\begin{proposition}[Frontier invariant]
After processing any prefix of the chosen crossing order, $F_i(\pi)$ is exactly
the sum of the corresponding partial state weights. At the final stage,
$F_n(\varnothing)=B_D(a)\pmod m$.
\end{proposition}
\begin{proof}
Before any crossing, the empty partial state has weight one. Every partial state
at stage $i+1$ is uniquely obtained by extending a state at stage $i$ with one of
two smoothings. The transition multiplies by exactly its smoothing coefficient
and the factors for newly closed circles. Summing states with identical outgoing
matchings is distributivity in the coefficient ring. Induction proves the first
claim. There are no boundary ends after the last crossing, giving the second.
The empty-PD single-circle case is handled separately with value $\delta$.
\end{proof}

\subsection{Resource and width bounds}
Let $w$ be the largest number of frontier ends. A safe state-count bound is
\begin{equation}
S(w)\le(w-1)!!=2^{O(w\log(w+2))}\qquad(w>0\text{ even}),
\label{eq:matchingbound}
\end{equation}
with $S(0)=1$. It is tempting to replace this by a Catalan number, but that
requires an appropriate common cyclic ordering on a disk boundary. An arbitrary
partial crossing scan can have several boundary components. We do not use an
unproved disk-boundary assumption in the complexity claim.

A deliberately conservative upper bound for one transition in the current
Python representation is $O((w+1)^2)$ elementary operations, including the
unranked disjoint-set operations and canonical sorting. For fixed-size modular
arithmetic, the uncapped evaluation therefore takes
\[
O\bigl(nS(w)(w+1)^2\bigr)
\]
operations, in addition to finding an order, and stores $O(S(w)(w+1))$ label and
residue words. For variable $m$, include the cost of arithmetic on $O(\log m)$-bit
integers and the initial modular inversions.

The \emph{default} screen stops as soon as the partially built dictionary exceeds
4096 entries. At most 4097 entries are present at the triggering instant, though
the previous dictionary also remains live. The cap can trigger before later
cancellations would shrink that dictionary; this sacrifices an opportunity to
screen, not correctness. Hitting this local state ceiling skips the screen and
continues to Alexander/Khovanov. Exhausting the global time budget gives
\Unknown. Thus an unlucky Jones evaluation cannot become an unbounded mandatory
exponential prepass in the default recognizer.

\section{Visible sums and compressed homology}
\label{sec:factor}
\subsection{Detecting and checking a separating circle}
For the connected plane projection graph $G$, compute the incident faces of
each edge. Two edges with the same two \emph{distinct} incident faces are
parallel edges in the plane dual. Their union in the dual supplies a two-edge
cycle and hence a candidate circle meeting $G$ in exactly two edge interiors.

The implementation groups edges by their unordered pair of incident face IDs.
It tests consecutive entries in each group; any two distinct parallel dual edges
already give a dual 2-cycle, so their order in the group need not be geometric.
For a proposed pair $\{e,f\}$, a graph traversal after deleting those two edges
finds a proposed vertex side $V_1$. Before accepting the cut, the program checks
that $V_1$ is nonempty and proper, that its entire edge boundary is precisely
$\{e,f\}$, that the two shared dual endpoints are distinct, and that both sides
are connected. It closes each side by identifying the two cut labels, and checks
both resulting PDs for spherical rotation systems and exactly one strand component.

\begin{lemma}[Meaning of an accepted cut]
An accepted cut decomposes the represented knot as the connected sum of the two
returned knot diagrams, with no crossing added or removed.
\end{lemma}
\begin{proof}
The two face arcs joining the cut-edge interiors form an embedded dual 2-cycle.
It is a Jordan curve in the projection sphere and avoids all crossings. The
checked edge boundary means it meets the projection exactly twice. A narrow
three-dimensional thickening gives the usual separating sphere meeting the knot
in two points. Each side is a one-strand tangle. Capping the two endpoints along
the separating circle is exactly the label identification performed by the code.
Its inverse operation is connected sum. The crossings and their cyclic port data
are inherited unchanged on each side.
\end{proof}

All accepted cuts are recorded in a tree with parent and child node IDs, the two
cut-edge labels, and the crossing indices on one side. The verifier rebuilds and
checks each cut, then validates the leaf list. An iterative worklist avoids the
Python recursion limit. A conservative bound for the current repeated-search
implementation is $O(n^3)$ word operations: $O(n)$ candidates, each checked in
$O(n)$ time, at each of at most $n-1$ splits. Sorting and integer-label bit costs
add polynomial overhead. This is not an optimized general graph-bond routine.

This decomposition is \emph{not topological prime factorization}. A composite
knot can have a projection with no visible two-edge sum cut. No conclusion about
primality is made when the procedure returns one factor.

\subsection{Reduced ranks and the complete detector}
\label{sec:rank}
Write
\[
r(K)=\dim_{\F}\rKh(K;\F),\qquad
P_D(t)=\sum_h\dim_{\F}\rKh^h(D;\F)\,t^h.
\]
Here the code records the \emph{raw cube degree} $h$, the number of 1-smoothings,
not the usual writhe-shifted homological degree, and does not compute quantum
gradings. Raw degrees are additive under cuts that preserve the set of crossings.

Over $\F$, unreduced homology consists of two copies of reduced homology with
quantum shifts and no homological shift \cite[Corollary 3.2.C]{Shumakovitch}.
Consequently every unreduced homological-degree rank used here is even.
An algebraic explanation is useful. On a resolution with circle variables
$x_1,\ldots,x_c$, let $X$ multiply by the basepoint-circle variable and let
$\nu=\sum_i\partial/\partial x_i$. In characteristic two,
\[
X^2=\nu^2=0,\qquad X\nu+\nu X=\id.
\]
The maps commute with the Khovanov differential: this is checked on the
multiplication and comultiplication of $\F[x]/(x^2)$. Thus $X\nu$ is a chain
projection onto $\operatorname{im}X$, the reduced subcomplex up to quantum shift,
and the other summand is isomorphic to it via $X$ and $\nu$. Forgetting quantum
degree gives the asserted doubling.

The complete decision rule is
\begin{equation}
K\text{ is the unknot}\quad\Longleftrightarrow\quad r(K)=1.
\label{eq:detector}
\end{equation}
Universal coefficients relate the field to the standard detection theorem:
\[\dim_{\mathbb Q}\rKh(K;\mathbb Q)\le r(K).\]
Reduced rational homology is
nonzero, since its graded Euler characteristic is the normalized Jones
polynomial and that polynomial has value one at one. Hence $r(K)=1$ implies
reduced rational rank one, to which Kronheimer--Mrowka's theorem applies
\cite{KM}. Conversely, the unknot has reduced rank one over $\F$. We do not
assume that torsion-free rational computations can always replace the $\F$
calculation.

\begin{proposition}[Connected-sum product]
For the two diagrams returned by an accepted cut,
\[
P_D(t)=P_{D_1}(t)P_{D_2}(t),\qquad r(K_1\#K_2)=r(K_1)r(K_2).
\]
\end{proposition}
\begin{proof}
Place each basepoint on the arc being joined. A pair of smoothing states on the
two sides determines exactly one state of the sum. The two distinguished circles
are joined to one distinguished circle; in the reduced complex each carries the
fixed basepoint label $x$. All remaining circle labels are independent. This
identifies the reduced cube for the sum with the tensor product of the reduced
cubes, with additive cube degree. Saddle maps in either block agree under this
identification. The tensor-product differential has no sign issue over $\F$.
K\"unneth over a field yields the homological-degree convolution and then the
rank product.
\end{proof}

Every factor has positive reduced rank. Combining the product with
\eqref{eq:detector}, the sum is the unknot precisely when all factors are.
This justifies both the early nontrivial-factor exit and the all-factors-unknot
branch. No unimplemented prime-decomposition oracle is needed for this rule.

\subsection{Avoiding a materialized tensor product}
The raw scanner still produces one object per surviving unreduced generator.
The new \path{factored_khovanov_rank} instead computes each factor's reduced
homological polynomial and multiplies these polynomials using integer
coefficient convolutions. The output is a map from degree to a binary integer
rank, not a list of generators. Identical \emph{normalized PD tuples} can reuse
a factor computation. This is exact tuple memoization, not isotopy testing and
not a complete canonical representation of knots.

The sum of factor crossing counts is $n$, and all output degrees are between
zero and $n$. Coefficients have $O(n)$ bits under the elementary exponential
bound on cube size. Ordinary coefficient convolutions therefore have polynomial
bit complexity. Even without the memoization, repeated bounded-size factors cost
only polynomial time to process. The extra unreduced factor is inserted once:
\[
\dim_{\F}\Kh(K_1\#\cdots\#K_k;\F)=2\prod_{i=1}^k r(K_i),
\]
not $\prod_i 2r(K_i)$.

\section{A faster exact scanning backend}
\label{sec:scan}
\subsection{The local algebra}
The baseline represents a boundary object by a matching. Morphisms are
$\F$-linear combinations of dotted cobordisms in a normal form indexed by dot
subsets on circles obtained by closing source against target. Composition uses
the Frobenius algebra $A=\F[x]/(x^2)$, with
\[
\begin{array}{c|cccc}
&1\otimes1&1\otimes x&x\otimes1&x\otimes x\\\hline
m&1&x&x&0
\end{array},\qquad
\Delta(1)=1\otimes x+x\otimes1,\quad\Delta(x)=x\otimes x.
\]
An undotted sphere is zero, a once-dotted sphere is one, and an extra dot gives
zero. A handle gives $m\Delta=0$ in characteristic two. Delooping a circle
replaces it by two summands. These are the same mathematical operations as in
the supplied backend, not a different homology theory \cite{BNtangles,BNfast}.

For an object consisting of $r$ arcs, the endomorphism algebra in this normal
form is
\begin{equation}
R_r=\F[x_1,\ldots,x_r]/(x_1^2,\ldots,x_r^2).
\label{eq:localring}
\end{equation}
A squarefree monomial is stored as a set of dotted components, and a polynomial
as the set of monomials with coefficient one. Addition is symmetric difference.

\begin{lemma}[Every unit is self-inverse]
An element of $R_r$ is a unit exactly when its constant coefficient is one.
Every such unit satisfies $u^{-1}=u$.
\end{lemma}
\begin{proof}
Write $u=\epsilon+\eta$, where $\eta$ has zero constant coefficient. In
characteristic two, commutativity gives $\eta^2=\sum_M M^2=0$, since every
nonconstant squarefree monomial has a squared variable. If $\epsilon=1$,
$u^2=1$; if $\epsilon=0$, $u^2=0$, so $u$ cannot be invertible.
\end{proof}

The baseline builds an inverse by a finite nilpotent series using general
surface composition. The replacement needs only a copy of the input monomial
set. It also recognizes identity compositions, and multiplies endomorphism
monomials directly: intersecting dot sets give zero; otherwise their union is
the product. General source/intermediate/target combinations still use the full
surface calculation.

\subsection{Separating topology from decorations}
The topological part of composing $a\to b\to c$ depends only on the three
matchings. It determines the connected surface components, their Euler
characteristics, and their outgoing boundary circles. Dot decorations are then
distributed over this fixed plan. Previously this topology was repeatedly
reconstructed for different decorations.

The implementation now caches topology plans, complete decorated compositions,
and the local crossing-tensor/delooping entries. Cached decorated values are
immutable; public morphism results are fresh mutable sets, so subsequent
symmetric differences cannot corrupt cached results. The entry ceilings are
4096 each for circle decompositions, gluing records, and composition plans;
16384 for decorated compositions; and 8192 for crossing entries.
\path{clear_caches()} makes a cold run explicit, and \path{cache_info()} exposes
hits and misses.

These are bounds on \emph{cache entry counts}, not a bound on total bytes.
A single morphism can have exponentially many monomials. The underlying
low-level topology records are internal shared objects and must be treated as
read-only; the tests specifically check nonmutation of returned public morphisms.
No polynomial-space theorem follows merely from adding an LRU limit.

\subsection{Cancellation and fill-in control}
If a differential component $\phi:b\to c$ is invertible, Gaussian cancellation
removes $b,c$. For any remaining incoming component $\delta:a\to c$ and
outgoing component $\gamma:b\to f$, replace the surviving differential by
\begin{equation}
d'_{af}=d_{af}+\gamma\phi^{-1}\delta.
\label{eq:schur}
\end{equation}
The sign is plus over $\F$. This is the Schur complement form of
Bar--Natan's Gaussian elimination lemma and preserves chain homotopy type
\cite{BNfast}. Algebraically, triangular basis changes isolate the isomorphism
$\phi$ as a contractible two-object summand; deleting that summand gives
\eqref{eq:schur}. The adjacent differentials transform consistently because
$d^2=0$.

Correctness does not prescribe the order of unit cancellations. Performance can
change sharply because different orders create different intermediate entries.
For a candidate $b\to c$, define
\[
q(b,c)=(|\operatorname{inc}(c)|-1)(|\operatorname{out}(b)|-1).
\]
It estimates the number of possible Schur corrections, not their actual
nonzero count. A heap prioritizes small $q$, breaking ties by the sum of the
two degrees and then object IDs. Removed or nonunit candidates are discarded.
An underestimated stale priority is refreshed and reinserted; a stale
\emph{overestimate} may be accepted. Thus \texttt{minfill} is a lazy
fill-in-aware heuristic, \emph{not} a guarantee of the globally cheapest pivot
at every step. Every accepted cancellation is nevertheless exact.

The old last-in-first-out choice remains available as
\path{pivot_strategy="stack"} in the raw Python API. When there are no outgoing
correction targets, the implementation skips even the unnecessary incoming
half-compositions. Otherwise $\phi^{-1}\delta$ is computed once per incoming
entry and reused. Composition statistics in this version count these actual
requests; the old counter charged some half-compositions differently. Timings
are compared across versions, but the old and new raw composition counters are
not presented as identical performance units.

\subsection{Faster crossing orders}
The baseline's greedy order repeatedly examines all remaining crossings.
The priority first maximizes the number of edges shared with the current
frontier, then prefers self-loops, then the lowest crossing index. For an
unprocessed crossing, the shared-edge count changes only when one of its
neighbors is processed. It never decreases before that crossing is selected.
The replacement builds adjacency counts once and maintains a lazy heap of these
same priorities.

\begin{proposition}[Order construction]
For any validated $n$-crossing PD and specified starting crossing, the new
routine returns exactly the baseline greedy order using $O(n\log(n+2))$
word operations and $O(n)$ words of storage.
\end{proposition}
\begin{proof}
There are $2n$ graph edges, counted with multiplicity. Each edge between different
crossings causes at most one priority update while its second endpoint remains
unprocessed. Initial heap entries and all updates therefore number $O(n)$.
Each heap operation costs $O(\log(n+2))$. An obsolete priority is ignored, and
the current priority of every eligible crossing is present. Lexicographic
comparison implements exactly the baseline tie rules. Loops contribute only
the fixed tie score and cause no neighbor update.
\end{proof}

Trying $r$ sampled starts costs $O(rn\log(n+2))$ instead of $O(rn^2)$.
The inherited start-sampling expression can use somewhat more than the requested
number of starts, but still $O(r)$; the default is a fixed number independent of
$n$. Asking for \emph{all} starts gives $O(n^2\log n)$, not the fixed-start bound.
The resulting order is heuristic: no optimal-cutwidth claim is made.

\section{Preprocessing, validation, and evidence}
\label{sec:safety}
\subsection{A linear descending test}
A knot traversal has $2n$ crossing visits. Relative to a chosen start, mark a
crossing bad when its first visit is under. Moving the start past one visit
changes the first-visit choice of exactly that crossing; all other first-visit
choices retain their order. Passing an over-visit increments the bad count,
and passing an under-visit decrements it. A zero count gives a descending start.
The routine performs an initial pass and a sliding pass in each traversal
direction, returning the least valid dart just as the baseline does. This gives
$O(n)$ word operations instead of testing $O(n)$ starts independently. A
descending diagram represents the unknot: lifting successive first encounters
above later ones permits the traversed arc to be straightened without changing
any crossing.

\subsection{An orientation fix needed for safe normalization}
The old sign routine looked only at the incoming over-port. A PD crossing may be
rotated by 180 degrees without changing it, and that operation can change the
absolute incoming port numbers. The corrected routine uses the relative
position of the incoming over-port $o$ and incoming under-port $u$:
\[
\text{sign}=+1\ \Longleftrightarrow\ (o-u)\bmod4=3.
\]
The relative difference is invariant under simultaneously rotating the ports by
two, and under reversing the strand orientation. The Alexander matrix was
updated at the same time to use the actual incoming and outgoing under-arcs;
changing only the sign routine would not be a safe patch. The bracket
normalization, mirroring, local port rotations, and Reidemeister tests use the
same convention throughout.

\subsection{Limits are not mathematical answers}
The raw scanning API now validates the PD and requires an explicit scan order to
be a permutation of \emph{all} crossing indices. Incomplete, repeated, negative,
out-of-range, and Boolean indices are rejected. Time limits must be finite and
nonnegative, and object and state ceilings must be positive integers when set.
The global deadline is shared across factor analyses and preprocessing.
Additional checks occur during scanning, simplification, and the Alexander
Bareiss calculation. New-object ceilings are tested while preparing the tensor
expansion, rather than only after allocating every new object.

Time limits are cooperative, not hard real-time or operating-system memory
limits. One sizeable elementary operation can overrun a deadline before the
next check. The object ceiling bounds active object creation, not morphism
bytes, caches, or all process memory. The benchmark harness adds an external
subprocess timeout for isolation. Under these limits, failure to complete returns
\Unknown; it must never be interpreted as \Knotted.

\subsection{Replayable evidence and its trust boundary}
Recognition outputs contain the Reidemeister trace, any descending start,
decomposition tree and factor results, modular evaluation data, Alexander
coefficients, or Khovanov ranks and scan order as appropriate. The
\path{verify_result} routine replays local moves and dual cuts; recomputes
Jones/Alexander obstructions; and, for a Khovanov verdict, recomputes the scan
with $d^2=0$ checks after every crossing. A knotted connected sum needs a verified
nontrivial factor; an unknot sum needs every factor verified as an unknot.

This is reproducible mathematical evidence, not a claim of a succinct
polynomial-size certificate for arbitrary Khovanov computations. The verifier
shares diagram and algebra routines with the recognizer. Its Khovanov branch
may be as expensive as computing the invariant again. Neither Python nor the
cobordism implementation has been formalized in Lean. Independent reference
algorithms and invariance tests reduce common-mode risk, but do not eliminate
that trust boundary.

\section{Complexity: achieved improvements and remaining barriers}
\label{sec:complexity}
\subsection{General worst case}
For an $n$-crossing connected diagram, there are $2^n$ resolutions and at most
$n+1$ circles in a resolution, so the unreduced full cube has at most
$2^{2n+1}$ generators. The analogous partial scanning objects, their dot-normal
forms, and differential entries have elementary exponential upper bounds as
well. Gaussian elimination and composition use a fixed number of polynomial
loops in these exponentially bounded quantities. Consequently an unbounded
execution of the supplied style of algorithm has the conservative bound
$2^{O(n)}$, including polynomial-size bit bookkeeping. The bounded Jones screen
and the polynomial diagram operations do not worsen this bound.

Neither the algebraic shortcuts nor the pivot heuristic changes that global
worst-case classification. In particular, no empirical regression on a few timing
points is used to infer an asymptotic exponent.

\subsection{Why width alone does not bound a materialized Khovanov complex}
\begin{theorem}[A fixed-width multiplicity obstruction]
There is a family of knot diagrams with $n=O(k)$ crossings and scan orders of
bounded frontier width for which the final reduced Khovanov rank is $r^k$ with
$r>1$. Any scanner that represents each final homology generator by a separate
object needs exponentially many final objects on this family.
\end{theorem}
\begin{proof}
Fix a diagram of a nontrivial knot $K$. Join $k$ copies in a visible linear
connected sum. The copies have bounded size and can be scanned one at a time;
only a bounded number of connection edges and one fixed-size block are active.
Thus there is a constant-width order, independent of $k$. By the connected-sum
formula, the reduced rank is $r(K)^k$. By \eqref{eq:detector}, $r(K)>1$.
An explicit final basis requires at least that many objects, independent of how
unit eliminations were scheduled earlier.
\end{proof}

The integer $r^k$ itself has only $O(k)$ bits. Thus this is a barrier to the
\emph{representation used by the supplied scanner}, not an output-size lower
bound for computing a rank in binary, and not a lower bound for all unknot
recognition algorithms. It explains why ordinary memoization of repeated
matching \emph{types} does not suffice: many distinct surviving generators can
have exactly the same matching type.

This also corrects a separate width statement in the supplied README. An
$n$-vertex planar graph of maximum degree four admits a vertex scan order with
$O(\sqrt n)$ crossing edges, rather than having optimal width $\Theta(n)$.
Indeed, recursively take balanced planar separators of $O(\sqrt n)$ vertices
\cite{LiptonTarjan}, process a separator and then its remaining components, and
sum the at-most-four incident edges per separator vertex along the active
recursion chain:
\[
O(\sqrt n)\sum_{j\ge0}(\sqrt{2/3})^j=O(\sqrt n).
\]
The implementation's greedy ordering need not find such an order, and we do not
implement a separator-order algorithm in this release. More importantly, even
\emph{constant} width does not remove the generator-multiplicity obstruction.
A boundary-only complexity assertion for the raw Khovanov representation would
therefore remain false after improving the order.

\subsection{An exponential improvement for the actual recognition pipeline}
\begin{theorem}[Structured-family speedup]
Fix a nontrivial Alexander-trivial knot $K$ and a fixed diagram for it. On the
family of visibly spliced copies $K^{\#k}$, the supplied baseline recognition
pipeline, if run to completion, uses $\Omega(r(K)^k)$ final Khovanov objects.
The improved recognizer and its factored rank routine run in polynomial time
in $k$ when the validated decomposition extracts bounded-size blocks.
\end{theorem}
\begin{proof}
The represented knot is nontrivial by the rank product, so sound Reidemeister and
descending tests cannot certify it as an unknot. Its Alexander polynomial is
one, because Alexander polynomials multiply under connected sum. The baseline
therefore reaches its Khovanov routine. Simplifying the diagram does not change
the final rank $r(K)^k$, so that routine must materialize exponentially many
objects.

In the improved pipeline, finding all validated visible cuts costs polynomial
time. Each extracted fixed-size block can be analyzed in constant time with
respect to $k$, even if it requires its complete Khovanov computation. There
are $O(k)$ blocks. Recognition may stop at the first certified nontrivial one.
For the factored rank routine, the remaining integer coefficient convolutions
have polynomial bit complexity. No tensor-product basis is expanded.
\end{proof}

The Conway fixture realizes the intended difficult-filter case: its Alexander
polynomial is one, while independent cube and scanning computations both give
$r=33$. The test family is generated by the included \path{connected_sum}
helper, and all decomposition cuts are replayed. This theorem is more than a
constant-factor micro-optimization, but its hypothesis is a visible decomposition,
not arbitrary topological compositeness.

\begin{corollary}[A precise restricted quasipolynomial bound]
Let $b$ be the largest crossing count of any block actually returned by the
validated diagram decomposition. The improved complete recognizer has the
conservative bound
\[
T(n,b)\le\poly(n)+n\,2^{O(b)}.
\]
In particular, bounded $b$ gives polynomial time; $b=O(\log n)$ also gives
polynomial time; and $b=O(\log^2 n)$ gives $n^{O(\log n)}$.
\end{corollary}
\begin{proof}
The factor crossing counts sum to $n$, hence there are at most $n$ nonempty
blocks. Each block's complete analysis has exponential upper bound in its own
crossing count, at most $2^{O(b)}$, plus polynomial preprocessing. Decomposition
and bookkeeping contribute the polynomial term. Substitution gives the three
specializations.
\end{proof}

When no visible cut exists, $b=n$ and this is only the ordinary exponential
bound. The corollary does not establish the requested global quasipolynomial
bound, nor does it say that every composite knot has a conveniently factored
projection.

\section{Reproducible measurements}
\label{sec:bench}
\subsection{Protocol and environment}
All numbers in Tables~\ref{tab:recognition}--\ref{tab:order} come from the
included \path{benchmarks/comparison.json} and are regenerated into LaTeX by
\path{tools/make_report_tables.py}. The measured environment was CPython
3.13.5, GCC 14.2.0, Linux x86\_64 with glibc 2.41. The processor reported
\texttt{AMD EPYC 9V74 80-Core Processor}; the container reported five logical
CPUs. Timings are single-process measurements in a shared container, not
isolated-hardware performance guarantees.

Each sample runs in a fresh child process, so caches start cold. The reported
timer excludes imports, process startup, JSON parsing, and initial
\path{Diagram} construction. Thus these are algorithm timings on an already
constructed diagram, not command-line wall times. The improved raw scan performs
its own internal revalidation, which \emph{is} counted; this gives it a small
extra safety cost absent from the old raw function. Every successful group has
three samples; the tables use the median. Raw minimum and maximum times and all
full result records are retained. A censored group is attempted once, with a
30-second cooperative limit and an outer subprocess timeout three seconds later.

The comparison is against the unchanged supplied source, not a recreated
approximation. Small-input timings are especially sensitive to runtime noise.
The last source hardening added only argument rejection and clarified comments;
valid benchmark paths are unchanged. The final 48-test suite was then rerun.

\subsection{Default recognition}
\begin{table}[htbp]
\centering\small
\begin{tabular}{lrrr}
\toprule
Input & Baseline (ms) & Improved (ms) & Ratio \\
\midrule
Conway (11 crossings) & 35.090 & 0.834 & $42.09\times$ \\
Kinoshita--Terasaka (11) & 37.594 & 0.904 & $41.58\times$ \\
Hard unknot (8) & 1.622 & 2.244 & $0.72\times$ \\
Trefoil (3) & 0.092 & 0.229 & $0.40\times$ \\
Figure-eight (4) & 0.125 & 0.272 & $0.46\times$ \\
Torus $T(3,5)$ (10) & 0.588 & 0.860 & $0.68\times$ \\
Grid-scrambled unknot & 0.320 & 0.280 & $1.14\times$ \\
Unknot braid (40) & 6.448 & 6.264 & $1.03\times$ \\
Four-strand braid (41) & 10.679 & 5.908 & $1.81\times$ \\
Five-strand braid (36) & 8.626 & 5.043 & $1.71\times$ \\
Conway $\#2$ (22) & 1,372.423 & 1.455 & $943.04\times$ \\
Conway $\#3$ (33) & $>30,000$ & 1.972 & $>15,210.39\times$ \\
\bottomrule
\end{tabular}

\caption{Default recognition, cold-process medians. Ratio is baseline time
 divided by improved time; ratios below one are regressions. The $>$ marker
 denotes a censored baseline, not a measured completion time.}
\label{tab:recognition}
\end{table}

The largest improvement among completed default comparisons is the two-Conway
sum, about $943\times$. The three-Conway baseline does not complete within the
30-second limit, while the new default completes in about 1.97\,ms. These
inputs are Alexander-trivial, which is precisely why the old pipeline computes
the large Khovanov complex. The single Conway and Kinoshita--Terasaka fixtures
are rejected quickly by the Jones screen without computing homology.

The result is not uniformly faster. The trefoil, figure-eight, and $T(3,5)$
fixtures were already rejected cheaply by Alexander; the added screen and cut
checks introduce overhead. The hard eight-crossing unknot is another regression
in default recognition: an agreeing Jones evaluation is correctly inconclusive
and the complete Khovanov fallback still has to run. Flags to disable Jones and
factorization permit the user to retain the old filter order when a workload
consists mainly of such easy cases.

\subsection{Raw scanning, separated from preprocessing}
\begin{table}[htbp]
\centering\small
\begin{tabular}{lrrr}
\toprule
Input & Baseline (s) & Improved (s) & Ratio \\
\midrule
Conway (11 crossings) & 0.034 & 0.015 & $2.30\times$ \\
Kinoshita--Terasaka (11) & 0.045 & 0.016 & $2.73\times$ \\
Hard unknot (8) & 0.002 & 0.002 & $1.11\times$ \\
Torus $T(3,5)$ (10) & 0.010 & 0.010 & $1.02\times$ \\
Four-strand braid (41) & 0.923 & 0.283 & $3.26\times$ \\
Five-strand braid (36) & $>30$ & 1.487 & $>20.18\times$ \\
Conway $\#2$ (22) & 1.395 & 0.235 & $5.93\times$ \\
\bottomrule
\end{tabular}

\caption{Raw Khovanov scans: no Reidemeister reduction, no Jones/Alexander
filter, and no connected-sum factorization. Both completed versions return the
same reduced ranks.}
\label{tab:scan}
\end{table}

This isolates improvement of the supplied backend itself. On the four-strand
41-crossing stress braid, the median drops from 0.923\,s to 0.283\,s; both
return reduced rank 109. On the two-Conway sum, the raw median drops from
1.395\,s to 0.235\,s. Both versions still materialize 2178 final unreduced
objects on that case, with 8118 objects at the pre-elimination peak. The gain
here is faster algebra, ordering, and cancellation management, not the compressed
sum algorithm.

For the five-strand 36-crossing fixture, the new raw scan returns reduced rank
2949 and unreduced rank 5898 in 1.487\,s median, with 17694 objects at the
pre-elimination peak and boundary width eight. The original raw scan does not
finish within our 30-second limit, so the comparison supports a conservative
speedup greater than $20\times$, not an exact ratio.

The archive's old \emph{historical} benchmark recorded a 600.147-second timeout
for this same raw case under Python 3.14.4 on Windows 11. That is useful for
identifying the stress input but is not a same-machine baseline. We do not divide
600 seconds by a Linux timing to advertise a speedup. Also, \emph{default}
recognition of the same 36-crossing input is already fast in the baseline:
8.626\,ms versus 5.043\,ms here. The dramatic raw-scan improvement must not be
confused with default end-to-end recognition on that input.

\subsection{Compressed rank computations}
\begin{table}[htbp]
\centering\small
\begin{tabular}{lrrrr}
\toprule
Input & Crossings & Reduced rank & Time (ms) & Peak factor objects \\
\midrule
Conway $\#2$ & 22 & $33^{2}$ & 24.783 & 246 \\
Conway $\#3$ & 33 & $33^{3}$ & 25.451 & 246 \\
Conway $\#8$ & 88 & $33^{8}$ & 30.007 & 246 \\
Conway $\#16$ & 176 & $33^{16}$ & 47.949 & 246 \\
Trefoil $\#8$ & 24 & $3^{8}$ & 3.374 & 18 \\
Trefoil $\#16$ & 48 & $3^{16}$ & 7.490 & 18 \\
\bottomrule
\end{tabular}

\caption{Complete reduced-rank and homological-degree computations using the
validated factorization. ``Peak factor objects'' is the largest pre-elimination
object count of any individual factor scan, not total process memory.}
\label{tab:factored}
\end{table}

For 16 Conway copies, the program returns
\[
33^{16}=1\,977\,985\,201\,462\,558\,877\,934\,081
\]
as reduced rank, together with the full raw homological-degree rank map, in
about 47.95\,ms median. Only two distinct normalized factor PDs are computed
in this fixture; the other 14 reuse an identical tuple. The maximal individual
factor scan has 246 objects before elimination. This does not mean the whole
process uses 246 objects' worth of memory: the decomposition records, cached
algebra, factor results, and integer polynomial maps are also stored.

On two Conway copies, the factored complete rank takes 24.78\,ms, as compared
with 235.42\,ms for the improved raw scan and 1395.42\,ms for the old raw scan.
This is distinct from the much cheaper recognition-only time, which stops after
one nontrivial factor is found.

\subsection{Ordering and pivot ablations}
\begin{table}[htbp]
\centering\small
\begin{tabular}{rrrr}
\toprule
$n$ & Baseline (ms) & Heap order (ms) & Ratio \\
\midrule
128 & 55.575 & 3.078 & $18.06\times$ \\
256 & 234.410 & 5.771 & $40.62\times$ \\
512 & 890.923 & 12.638 & $70.50\times$ \\
1024 & 3,742.642 & 81.868 & $45.72\times$ \\
\bottomrule
\end{tabular}

\caption{Order construction alone on the closure of
$\sigma_1\cdots\sigma_n$. Both algorithms return the same chosen order.
These timings include all sampled greedy starts, but not a homology scan.}
\label{tab:order}
\end{table}

The order-only experiment isolates the asymptotic improvement from
$O(n^2)$ to $O(n\log n)$ per start. On the $n=512$ fixture, the medians are
891\,ms and 12.64\,ms. At $n=1024$ the improved median is 81.87\,ms;
the observed curve is not asserted to be a precise asymptotic fit. Both
versions return the same orders in the differential order tests. On the
256-crossing unknot chain, the entire raw scan decreases from 250.04\,ms
to 22.80\,ms.

\begin{table}[htbp]
\centering\small
\begin{tabular}{lrr}
\toprule
Input & Cached algebra, stack (s) & Cached algebra, heap (s) \\
\midrule
Conway (11 crossings) & 0.014763 & 0.013504 \\
Kinoshita--Terasaka (11) & 0.015803 & 0.015057 \\
Five-strand braid (36) & $>10$ & 1.497090 \\
\bottomrule
\end{tabular}

\caption{Pivot ablation: both columns use the new cached algebra and order
routine. Fresh processes, median of three successful samples, ten-second
limit for the censored stack case.}
\label{tab:ablation}
\end{table}

Table~\ref{tab:ablation} is stored separately in
\path{benchmarks/ablation.json}. The small named examples gain only modestly
from changing stack to heap pivots once the algebra is already accelerated.
The 36-crossing case remains incomplete at ten seconds with stack pivots, while
heap pivots complete in 1.497\,s median. This supports the practical importance
of pivot order on that case; it is neither a general speed guarantee nor a proof
that the heuristic minimizes fill-in.

\section{Validation actually performed}
\label{sec:tests}
The final command
\begin{lstlisting}
PYTHONPATH=fast python -m unittest discover -s tests -v
\end{lstlisting}
completed with \textbf{48 tests passing in 8.538 seconds}; the full log is in
\path{benchmarks/tests.log}. These are test methods, several of which contain
many independently generated cases. The unchanged baseline's 19 tests were
also run separately and passed. The new test suite retains those 19 legacy
methods, explicitly disabling the new filters where an old method choice is
being asserted, plus 29 additional methods.

The algebra tests compare every one of the 128 units in a three-arc
endomorphism algebra against the old inverse routine. They compare 1200
random valid planar six-endpoint cobordism compositions against the unchanged
composition routine, and check that mutating a public result does not mutate
its cached value. Sixty random scans compare both pivot strategies and the old
backend, with differential-square checks. An independent dense reduced
cube-of-resolutions implementation is tested on 80 random small braid diagrams
and on the Conway, Kinoshita--Terasaka, hard-unknot, and $T(3,5)$ fixtures.

The Jones tests evaluate 100 random diagrams by an independent full state sum
on $4n$ crossing darts, at three choices of modular parameters. That reference
does not call the frontier transition or cobordism gluing code. Additional
checks cover small unknots, diagram relabeling and crossing reordering,
180-degree port rotations, mirrors, Reidemeister-II insertion and the braid
relation, nonprime moduli, and $\delta=0$. Equality is explicitly tested to remain
inconclusive.

Connected-sum tests compare raw and factored ranks and their degree
convolutions, mixed summands, repeated Conway factors, and sums of unknots.
Malformed cuts, duplicate tree data, and tampered evidence are rejected. The
linear descending test matches the baseline on 300 random diagrams; greedy
orders match on 90 random diagrams at every possible starting crossing.
Invalid orders, parameter values, ceiling exhaustion, skipped Jones screens,
CLI exit codes, and result replay are tested too.

The 36-crossing stress scan is tested with $d^2=0$ after every stage. A separate
script, \path{tools/validate_stress.py}, compares it with the diagram obtained
by six valid R2 moves, reducing 36 crossings to 24. Both complete scans return
reduced rank 2949, and their full raw homological-degree maps agree after the
expected shift by six. The reduced diagram's Alexander coefficient list is
\[
(2,-19,80,-206,372,-512,567,-512,372,-206,80,-19,2),
\]
whose absolute evaluation at $-1$ is also 2949. This determinant agreement is
a sanity check, \emph{not} a general theorem equating determinant and Khovanov
rank. The large-diagram invariance check uses the same scan backend on two
Reidemeister-related inputs, so it is not an independent large Khovanov
calculation. Independent dense-cube checks are limited to feasible small cases.

All completed baseline/improved benchmark pairs have consistent mathematical
outputs. Four included example recognition records have been replayed:
Conway, the hard eight-crossing unknot, three Conway copies, and the 36-crossing
braid. Their evidence files are in \path{examples/evidence}. These experiments
support correctness and reproduce the reported improvements; they do not prove
that no implementation bug remains on arbitrary inputs.

\section{Using and extending the implementation}
\label{sec:usage}
From the unpacked archive root, these commands need no installation:
\begin{lstlisting}
cd fast
python -m fastunknot recognize examples/conway.json
python -m fastunknot recognize examples/conway.json --output result.json
python -m fastunknot verify examples/conway.json result.json
python -m fastunknot jones examples/conway.json
python -m fastunknot khovanov examples/conway.json
python -m fastunknot khovanov examples/conway.json --no-factor
\end{lstlisting}
The last two commands distinguish compressed-factor homology from a raw scan.
To use the original pipeline stages with the improved backend:
\begin{lstlisting}
python -m fastunknot recognize examples/conway.json \
  --no-jones --no-factor
\end{lstlisting}
Optional flags include \path{--seconds}, \path{--max-objects},
\path{--jones-max-states}, \path{--check-d2}, and switches disabling individual
preprocessors. Exit code 0 means a completed mathematical result; code 2 means
invalid input or rejected evidence; code 3 means resource-limited \Unknown.
A standalone \texttt{jones} command returns an evaluation, not an unknot verdict.

For an importable installation, Python 3.10 or later is required:
\begin{lstlisting}
python -m pip install ./fast
\end{lstlisting}
There are no runtime dependencies beyond the standard library. Example Python use:
\begin{lstlisting}[language=Python]
from fastunknot import Diagram, recognize, verify_result
from fastunknot import khovanov_rank, factored_khovanov_rank

D = Diagram.from_braid(2, [1, 1, 1])
result = recognize(D, seconds=10)
print(result.status, result.method)
if result.status != "UNKNOWN":
    assert verify_result(D, result.to_json(), seconds=10)

raw = khovanov_rank(D.pd, pivot_strategy="minfill")
compressed = factored_khovanov_rank(D)
assert raw["reduced_rank"] == compressed["reduced_rank"]
\end{lstlisting}

The CLI keeps the existing input formats. For example,
\begin{lstlisting}
{"braid": {"strands": 2, "word": [1, 1, 1]}}
\end{lstlisting}
is a three-crossing braid closure. Arbitrary integer PD edge labels are
normalized by \path{Diagram.from_pd}. Direct construction of the low-level
\path{Diagram} dataclass bypasses those checks and is not the supported way to
validate untrusted input.

To reproduce the experiments and rebuild this report, return to the archive root:
\begin{lstlisting}
python tools/make_inputs.py
PYTHONPATH=fast python -m unittest discover -s tests -v
python tools/benchmark.py --repeats 3 --timeout 30
python tools/ablate.py --repeats 3 --timeout 10
python tools/validate_stress.py
python tools/make_report_tables.py
cd report
pdflatex -interaction=nonstopmode -halt-on-error report.tex
pdflatex -interaction=nonstopmode -halt-on-error report.tex
\end{lstlisting}
The first command regenerates the structured connected-sum and chain fixtures.
The original random-braid stress fixtures are already supplied and are preserved
verbatim in meaning; their source words and provenance are recorded. A quick
benchmark subset is available with \path{--quick}. On Windows, installation via
\texttt{pip} avoids shell-specific \texttt{PYTHONPATH} syntax. The test files
also set the import path themselves.

Useful further engineering work would include better exact workload scheduling
between Jones and Alexander, a disk-aware separator decomposition for the
Jones state table, and more general compression of surviving Khovanov
multiplicities. A tuple cache alone cannot provide the latter: interactions
between equal matching types may require nontrivial differential structure.
For the hierarchy route, the next deliverable should be independently checked
compressed cutting and boundary-pattern operations with explicit bit-size
contracts, not a renamed call to the exponential fallback.

\section{Conclusion}
The delivered recognizer preserves exact decision semantics and improves both
its early obstruction stages and its complete homological backend. The
characteristic-two inverse shortcut and lazy pivot strategy speed up genuinely
raw scans. Validated sum decomposition eliminates an exponential
representation cost on a concrete infinite family, including Alexander-trivial
examples that force the original recognition pipeline into Khovanov homology.
The factored rank routine demonstrates that a huge rank need not imply a huge
explicit basis.

The general $n^{O(\log n)}$ target remains unproved for this implementation.
The supplied Lackenby paper supplies real compressed-topology ingredients, but
turning the talk's logarithmic-depth strategy into an executable with complete
size and time accounting is a different construction. The accompanying source,
measurements, proofs, and test logs distinguish the improvements achieved here
from that remaining objective.

\clearpage
\begin{thebibliography}{10}
\raggedright
\bibitem{Input}
\emph{Knots.zip}, archive supplied with the task, specifically
\path{fast/fastunknot}, \path{fast/tests}, \path{fast/results/benchmark.json}, and
\path{docs}. The original Python sources and relevant benchmark provenance are
retained in this deliverable. Accessed 18 September 2026.

\bibitem{Ltalk}
Marc Lackenby, \emph{Unknot recognition in quasi-polynomial time}, February 2021.
Supplied \path{docs/quasipolynomial-talk.pdf}; 109 physical pages. In particular,
physical slides 14, 66--74, and 109.

\bibitem{L2026}
Marc Lackenby, \emph{Incompressible surfaces, hierarchies and unknot recognition},
arXiv:2607.23350v1, 2026. Supplied source
\path{docs/arXiv-2607.23350v1/algorithm-incompressible-250726.tex}.
\url{https://arxiv.org/abs/2607.23350}.

\bibitem{BNfast}
Dror Bar--Natan, \emph{Fast Khovanov homology computations},
Journal of Knot Theory and Its Ramifications \textbf{16} (2007), 243--255.
arXiv:math/0606318.
\url{https://arxiv.org/abs/math/0606318}.

\bibitem{BNtangles}
Dror Bar--Natan, \emph{Khovanov's homology for tangles and cobordisms},
Geometry \& Topology \textbf{9} (2005), 1465--1499.
\url{https://www.math.utoronto.ca/drorbn/papers/Cobordism/Cobordism.pdf}.

\bibitem{KM}
Peter B. Kronheimer and Tomasz S. Mrowka,
\emph{Khovanov homology is an unknot-detector},
Publications Math\'ematiques de l'IH\'ES \textbf{113} (2011), 97--208.
arXiv:1005.4346.
\url{https://arxiv.org/abs/1005.4346}.

\bibitem{Shumakovitch}
Alexander N. Shumakovitch, \emph{Torsion of the Khovanov homology},
Fundamenta Mathematicae \textbf{225} (2014), 343--364;
arXiv:math/0405474, v2 (2018), especially Corollary 3.2.C.
\url{https://arxiv.org/abs/math/0405474}.

\bibitem{Kauffman}
Louis H. Kauffman, \emph{State models and the Jones polynomial},
Topology \textbf{26} (1987), no.~3, 395--407.
\url{https://doi.org/10.1016/0040-9383(87)90009-7}.

\bibitem{LiptonTarjan}
Richard J. Lipton and Robert E. Tarjan, \emph{A separator theorem for planar graphs},
SIAM Journal on Applied Mathematics \textbf{36} (1979), no.~2, 177--189.
\url{https://www.jstor.org/stable/2100927}.
\end{thebibliography}
\end{document}
